% topic2.tex -- Module 8 Topic 2: Structure of the atom, Q15-28.
\topictitle{Topic 2 - Structure of the Atom}

% ---- Q15 (what Thomson showed)
\begin{quizq}{}
What did Thomson's experiments with cathode rays show?
\opts{Cathode rays are electromagnetic waves.}{Cathode rays are positively charged ions.}
  {Cathode rays are neutral atoms.}
  {Cathode rays are negatively charged particles, much lighter than any atom.}
\end{quizq}

% ---- Q16 (Thomson: speed from crossed fields)
\Needspace{13\baselineskip}
\begin{quizq}{}
In Thomson's experiment, cathode rays pass undeflected through an electric field of
\qty{2.0e4}{V.m^{-1}} and a perpendicular magnetic field of \qty{1.0e-3}{T}. What is the speed of the
cathode rays?
\opts{\qty{20}{m.s^{-1}}}{\qty{2.0e4}{m.s^{-1}}}{\qty{2.0e7}{m.s^{-1}}}{\qty{2.0e10}{m.s^{-1}}}
\end{quizq}

% ---- Q17 (Thomson: charge-to-mass ratio)
\Needspace{13\baselineskip}
\begin{quizq}{}
The electric field in the previous question is switched off. The cathode rays now move in a circle of
radius \qty{0.114}{m} in the magnetic field of \qty{1.0e-3}{T}. What charge-to-mass ratio does this give?
\opts{\qty{3.5e18}{C.kg^{-1}}}{\qty{1.8e11}{C.kg^{-1}}}{\qty{8.8e3}{C.kg^{-1}}}{\qty{5.7e-12}{C.kg^{-1}}}
\end{quizq}

% ---- Q18 (Millikan: number of excess electrons)
\Needspace{14\baselineskip}
\begin{quizq}{}
In a Millikan experiment, an oil drop of mass \qty{4.9e-15}{kg} is held stationary between two horizontal
plates \qty{1.5}{cm} apart, with a potential difference of \qty{450}{V} between them. How many excess
electrons does the drop carry?
\opts{1}{10}{100}{1000}
\end{quizq}

% ---- Q19 (Millikan's conclusion)
\Needspace{15\baselineskip}
\begin{quizq}{}
What was the main conclusion of Millikan's oil-drop experiment?
\opts{The charge on every drop was a whole-number multiple of about \qty{1.6e-19}{C}, so charge is
   quantised.}
  {Electrons have a mass of \qty{9.1e-31}{kg}.}{Oil drops are always positively charged.}
  {The charge on a drop can have any value.}
\end{quizq}

% ---- Q20 (Geiger-Marsden: plum-pudding prediction)
\Needspace{15\baselineskip}
\begin{quizq}{}
According to Thomson's plum-pudding model, what should have happened to the alpha particles that Geiger
and Marsden fired at a thin gold foil?
\opts{Most of them would bounce straight back.}{They would all be absorbed by the foil.}
  {Many would be deflected through large angles by the electrons.}
  {Almost all would pass through, with only very small deflections.}
\end{quizq}

% ---- Q21 (Rutherford: conclusion from large deflections)
\Needspace{16\baselineskip}
\begin{quizq}{}
Which conclusion did Rutherford draw from the few alpha particles that were deflected through more
than \ang{90}?
\opts{The positive charge and almost all the mass of an atom are concentrated in a tiny nucleus.}
  {An atom is mostly empty space.}{Electrons orbit in fixed energy levels.}
  {Gold atoms are negatively charged.}
\end{quizq}

% ---- Q22 (Chadwick: the neutron)
\Needspace{16\baselineskip}
\begin{quizq}{}
In Chadwick's experiment, beryllium was bombarded with alpha particles. The radiation produced knocked
protons out of paraffin wax. What did Chadwick conclude this radiation was?
\opts{High-energy gamma rays}{A stream of electrons}
  {Neutral particles with about the same mass as a proton}{Alpha particles bouncing off the beryllium}
\end{quizq}

% ---- Q23 (why not gamma rays)
\Needspace{17\baselineskip}
\begin{quizq}{}
Why did Chadwick conclude that this radiation was NOT gamma radiation?
\opts{It was deflected by a magnetic field.}
  {Gamma photons would need far more energy than was available to knock protons out at the observed
   speeds, but a neutral particle with a mass similar to a proton's could do this easily.}
  {It travelled faster than light.}{It had a negative charge.}
\end{quizq}

% ---- Q24 (balanced fields, double speed)
\Needspace{15\baselineskip}
\begin{quizq}{}
Protons moving at speed $v$ pass undeflected through a region where an electric field of
\qty{5.0e3}{V.m^{-1}} is perpendicular to a magnetic field. A proton then enters the same region at a speed
of $2v$. What is the size of its acceleration as it enters?
\opts{0}{\qty{4.8e11}{m.s^{-2}}}{\qty{9.6e11}{m.s^{-2}}}{\qty{1.4e12}{m.s^{-2}}}
\end{quizq}

% ---- Q25 (cathode rays in an electric field)
\Needspace{14\baselineskip}
\begin{quizq}{}
A beam of cathode rays passes between two horizontal parallel plates. The top plate is positive and the
bottom plate is negative. In which direction is the beam deflected?
\opts{Upwards, towards the positive plate}{Downwards, towards the negative plate}{It is not deflected}
  {Into the page}
\end{quizq}

% ---- Q26 (number of neutrons)
\Needspace{12\baselineskip}
\begin{quizq}{}
How many neutrons are in a nucleus of uranium-238 (\nuc{238}{92}{U})?
\opts{330}{238}{146}{92}
\end{quizq}

% ---- Q27 (plum-pudding model)
\Needspace{15\baselineskip}
\begin{quizq}{}
Which statement describes Thomson's plum-pudding model of the atom?
\opts{A tiny positive nucleus, with electrons orbiting around it}
  {Electrons in fixed energy levels around a nucleus}
  {A nucleus surrounded by a cloud showing where the electrons are likely to be}
  {A sphere of positive charge, with electrons embedded in it}
\end{quizq}

% ---- Q28 (Millikan: voltage doubled)
\Needspace{14\baselineskip}
\begin{quizq}{}
An oil drop is held stationary between two horizontal plates. The potential difference between the
plates is suddenly doubled. Ignoring air resistance, what happens to the drop?
\opts{It stays stationary.}{It moves downwards at a constant speed.}{It accelerates downwards.}
  {It accelerates upwards.}
\end{quizq}
